% Owner: @Capheuss (track A). Only the owner edits this file; send input to them on #40. Taint propagation and the Age Matrix.

% Issue #4 (A2): baseline STT taint propagation.
\paragraph{Taint propagation}
Taint is kept in two tables, one for integer and one for floating-point
physical registers, each holding one YRoT per register, in the style of BOOM's
busy table. The tables sit at the boundary between rename and dispatch rather
than inside the rename stage, because BOOM assigns each micro-op its ROB index
in the core, after rename. An instruction's YRoT is the youngest valid taint
among all of its register sources, including store data and the third operand
of fused multiply-adds; integer \texttt{x0} and immediates contribute nothing.
A load or atomic then writes its own ROB index to its destination, becoming a
new root, while any other instruction writes its YRoT. A load's own YRoT
covers only its address, so a load with an untainted address may still execute
under an unresolved branch. Instructions dispatched in the same cycle must see
each other's taint, so a bypass passes each lane's destination taint to the
younger lanes that read it. Taint tables are written only for lanes that
actually dispatch, so a stalled lane creates no root and recomputes its taint
when it dispatches later.

\paragraph{Visibility point and recovery}
The ROB keeps one pending bit per entry, set when a conditional branch or
\texttt{jalr} enters the ROB and cleared when the branch resolves or the entry
is killed, flushed or committed. The VP is the oldest entry whose bit is set,
found with BOOM's existing age-priority encoder from the ROB head, and is
marked invalid when no branch is pending. A branch that resolves in one cycle
releases its dependent transmitters in the next.

Every cycle, every table entry whose YRoT has become safe is cleared. This eager 
clearing keeps every tag valid as a ROB index. A root can only retire from the 
ROB head, where every unresolved branch is younger than it, so its tag is always 
cleared before its ROB slot can be reused. 
It also means a branch mispredict needs no snapshot of the taint tables. 
Surviving registers keep their taint and continue to clear, and a register freed 
by a squashed instruction returns to BOOM's free list, where dispatch overwrites 
its taint when the register is next allocated. A full flush clears both tables, since
committed values are untainted. During recovery, \texttt{stt\_block} holds
every transmitter. It is set for either phase of a mispredict, during a
rollback, and for the two cycles after a flush.

% Issue #5 (A3): Age Matrix.
\paragraph{The $\mu$STT Age Matrix}
In baseline STT, the same-cycle bypass chains from lane to lane so each lane's
YRoT can depend on the result of every older lane in the group, and its delay
grows with dispatch width. The \emph{Age Matrix} breaks this chain. It
compares the taint read from the table for every source in the group pairwise,
in parallel with the dependency check, so no comparison waits for an older
lane's result. Each lane then finds the older lanes it reaches through
same-cycle dependencies and selects the youngest table taint among their
sources. If a lane reaches a load in its own group, that load's own ROB index
wins without a comparison, since a same-group root is younger than any taint
already in the table. We port $\mu$STT's Verilog to Chisel at BOOM's widths,
behind \texttt{enableAgeMatrix}, sharing the tables, clearing and writes with
the bypass. A structural change simplifies the port from $\mu$STT. 
Instead of computing two separate dependency closures, ours computes a single one 
(reachability through lanes that are not loads) and selects exactly the same set of
taints. Building the bypass also exposed a scaling problem. Verilog
generation inlined each lane's taint into every younger lane, so the
expression grew exponentially with width, and at width 4 Verilator could not
compile it. Keeping each lane's taint as a named wire fixes this without
changing the logic.
